\section{The physical origin of intelligence: active error correction and reverse compilation}
\label{sec:intelligence}

In a universe where phase friction is unavoidable for finite observers, stable low-entropy structure requires continuous repair. This section proposes an operational definition of life and intelligence as a phase transition: the emergence of predictive feedback loops that actively suppress local mismatch production while exporting dissipation.

\subsection{Agents as predictive-feedback subsystems}

\begin{definition}[Agent]
\label{def:agent}
A subsystem $\mathcal{S}$ is an \emph{agent} if, under finite computational flux $E$, it can:
\begin{enumerate}
  \item acquire conditional information about its environment through a finite-resolution readout instrument (O5),
  \item maintain an internal model that predicts future readout outcomes of the scan,
  \item implement feedback operations that reduce its own phase-friction entropy production rate relative to the passive baseline.
\end{enumerate}
\end{definition}

\noindent This definition converts Maxwell-demon language into engineering constraints: the demon does not violate the second law because it must pay readout and erasure costs \cite{Szilard1929,Landauer1961,Bennett1973}. What an agent accomplishes is not entropy elimination but entropy \emph{redirection}: moving dissipation into an external waste channel while protecting internal structure.

\subsection{Reverse compilation in the \texorpdfstring{$\Omega$}{Omega} dictionary}

In the $\Omega$ implementation picture, local update rules are exactly compilable into nearest-neighbor circuits; the required depth is the routing overhead $\kappa(x)$ (Section~\ref{sec:lapse_thermo}). ``Forward compilation'' maps dynamics to required overhead; ``reverse compilation'' is the agentic operation: using limited internal bits $I$ to infer which future phase information will be read out and to preconfigure control operations that cancel mismatch.

In HPA language, prediction corresponds to compressing future readout words in the canonical Zeckendorf/Ostrowski coordinate system, while control corresponds to applying phase corrections compatible with the Weyl structure. The net effect is a local reduction of mismatch density $\sigma(x)$ and therefore a reduction in $S_{\mathrm{pf}}$ production.

\subsection{Geometric Landauer principle}

Landauer's bound provides the minimal energetic cost of erasing one bit at temperature $T$ \cite{Landauer1961}. In HPA--$\Omega$, the operational temperature is $T_c$ and the erasure must be implemented through a geometric communication network with routing overhead. This motivates a refined principle:
\begin{equation}
W_{\mathrm{erase}} \ge k_B T_c \ln 2 + Z_{\mathrm{geom}},
\label{eq:geometric_landauer}
\end{equation}
where $Z_{\mathrm{geom}}$ is a geometric impedance term capturing the extra work required to route and re-encode information through constrained locality. Appendix~\ref{app:geometric_landauer} records a minimal formulation consistent with the additivity of log-costs.

\subsection{Information-thermodynamic bound on predictive gain}

Feedback control admits a quantitative upper bound: the work (or free-energy) advantage achievable by measurement-and-feedback is limited by the information acquired. In stochastic thermodynamics, generalized second-law inequalities take the form
\begin{equation}
  \langle W\rangle \ge \Delta F - k_B T\, I,
\end{equation}
where $I$ is an appropriate mutual-information term between measurement outcomes and system degrees of freedom \cite{SagawaUeda2010,ParrondoHorowitzSagawa2015}. Interpreting $T$ as the computational temperature $T_c$ in HPA--$\Omega$, we obtain an operational bound on predictive free-energy gain:
\begin{equation}
  \dot F_{\mathrm{pred}} \ \le\ k_B T_c\,\dot I_{\mathrm{pred}},
\label{eq:info_bound}
\end{equation}
where $\dot I_{\mathrm{pred}}$ denotes the mutual-information rate between the agent's internal state and future readout outcomes (a ``thermodynamics of prediction'' interface) \cite{StillSivakBellCrooks2012}.

\subsection{Survival criterion and predictive efficiency}

Let $\dot F_{\mathrm{pred}}$ denote the rate of free-energy gain achieved by prediction-and-control, and let $\dot W_{\mathrm{diss}}$ denote the dissipation rate required to maintain structure against phase friction. A necessary condition for sustained existence is
\begin{equation}
\dot F_{\mathrm{pred}} > \dot W_{\mathrm{diss}}.
\label{eq:survival}
\end{equation}
Combining Eq.~\eqref{eq:survival} with the information bound \eqref{eq:info_bound} yields a directly testable necessary condition in terms of an information rate:
\begin{equation}
  \dot I_{\mathrm{pred}} > \frac{\dot W_{\mathrm{diss}}}{k_B T_c}.
\end{equation}
Define the predictive efficiency
\begin{equation}
\eta_{\mathrm{pred}} := \frac{\dot F_{\mathrm{pred}}}{\dot W_{\mathrm{diss}}}.
\end{equation}
On this view, evolution is the search for architectures that maximize $\eta_{\mathrm{pred}}$ under constraints imposed by readout resolution, routing overhead, and mismatch density. Intelligence is then a phase transition in which an error-correcting feedback loop becomes self-sustaining and scalable in Zeckendorf depth.


